\section{Conclusion}
Current-cohort forecasting creates a practical bridge between temporal event
prediction and operational uncertainty. Across a cluster and a workflow,
holding individual probabilities fixed reveals substantial benefits from
shared variation, a smaller but measurable contribution from observed
relations, and no consistent advantage from richer context networks. These
benefits depend on the application and the evaluation target: better count
CRPS can coexist with worse reserve loss, and simple aggregate-error
resampling remains competitive. The benchmark and reference implementation
make these distinctions reproducible. Extending the task to arrivals,
repeated transitions, and additional organizations is needed before treating
such scenarios as a general operational simulator.
